% 14
\begin{frame}[t]{前六周的学习与复现安排}
\vspace{.15cm}
\centering
\begin{tikzpicture}[x=1.3cm,y=.78cm,
  rl/.style={anchor=east,font=\small,align=right},
  bar/.style={rounded corners=2pt,fill=PaleTeal,draw=DarkBlue,line width=.5pt},
  ms/.style={diamond,fill=DarkBlue,inner sep=2pt},
  ml/.style={font=\scriptsize,text=Muted,anchor=west,align=left}]
 \foreach \w in {1,...,6}{\node[font=\scriptsize\bfseries,text=DarkBlue] at (\w-.5,0) {第 \w\ 周};
   \draw[draw=Line,line width=.4pt] (\w-1,-.35) -- (\w-1,-4.4);}
 \draw[draw=Line,line width=.4pt] (6,-.35) -- (6,-4.4);
 \node[rl] at (-.2,-1) {\ico{\faBookOpen}\ 精读与任务定义};
 \draw[bar] (0,-1.3) rectangle (2,-.7);
 \node[ms] at (2,-1) {}; \node[ml] at (2.12,-1) {任务定义 + 评分脚本};
 \node[rl] at (-.2,-2) {\ico{\faRuler}\ 数据审计与基线};
 \draw[bar] (1,-2.3) rectangle (4,-1.7);
 \node[ms] at (4,-2) {}; \node[ml] at (4.12,-2) {基线表 + 可复现脚本};
 \node[rl] at (-.2,-3) {\ico{\faMicrochip}\ 接入公开预测器};
 \draw[bar] (3,-3.3) rectangle (6,-2.7);
 \node[ms] at (6,-3) {}; \node[ml] at (6.12,-3) {增益与分层结果};
 \node[rl] at (-.2,-4) {\ico{\faHospital}\ 并行：合作条件确认};
 \draw[bar,fill=PaleAmber,draw=Muted,dashed] (0,-4.3) rectangle (6,-3.7);
 \node[ml] at (6.12,-4) {伦理与数据流向方案};
\end{tikzpicture}\par
\vspace{.7cm}
\raggedright\textbf{阶段判断：}预训练模型是否超过基线？增益是否随背景与扰动类型变化？
\source{拟议安排，需结合课程、计算预算、数据规模与生物合作者反馈调整。}
\end{frame}
